% !TEX root = ../main.tex
% File: part1/chapters_reasoning/sec4_rl_reasoning.tex

\section{DeepSeek-R1: một mô hình suy luận được huấn luyện như thế nào}
\label{sec:rl_reasoning}

Chương~\ref{chap:post_training} đã dựng toàn bộ bộ máy: phần thưởng kiểm chứng được (mục~\ref{ssec:rlvr}), thuật toán GRPO (mục~\ref{ssec:grpo}) và họ biến thể sau nó (mục~\ref{ssec:grpo_variants}). Mục này dùng chúng để đọc một trường hợp cụ thể -- mô hình suy luận mở đầu tiên được mô tả đầy đủ công thức huấn luyện.

\subsection{DeepSeek-R1-Zero: suy luận “nảy sinh” từ học tăng cường thuần túy}
\label{ssec:r1_zero}
DeepSeek-R1 \cite{deepseekai2025r1} bắt đầu bằng một thí nghiệm: áp dụng GRPO \textbf{trực tiếp lên mô hình nền} DeepSeek-V3-Base, \emph{không có bất kỳ bước SFT nào}.

\paragraph{Thiết lập tối giản.}
\begin{itemize}
	\item \textbf{Mẫu prompt:} Chỉ yêu cầu về cấu trúc -- mô hình phải viết quá trình suy luận trong cặp thẻ \texttt{<think>}\ldots\texttt{</think>} rồi đưa ra câu trả lời trong \texttt{<answer>}\ldots\texttt{</answer>}. Không có gợi ý nào về \emph{cách} suy luận.
	\item \textbf{Phần thưởng chính xác (accuracy reward):} Đáp án toán ở định dạng quy định (ví dụ trong hộp) được so khớp tự động; bài lập trình được chấm bằng trình biên dịch và test case.
	\item \textbf{Phần thưởng định dạng (format reward):} Thưởng khi quá trình suy nghĩ nằm đúng trong thẻ \texttt{<think>}.
	\item \textbf{Không dùng reward model nơ-ron} (cả theo kết quả lẫn theo quá trình), vì lo ngại reward hacking ở quy mô lớn và chi phí huấn luyện lại reward model.
\end{itemize}

\paragraph{Kết quả.} Sau hàng nghìn bước RL, pass@1 trên AIME 2024 tăng từ \textbf{15.6\%} lên \textbf{71.0\%}; với bỏ phiếu đa số trên 64 mẫu đạt \textbf{86.7\%}, ngang mô hình o1 bản đầu của OpenAI.

\paragraph{Sự tự tiến hóa và “khoảnh khắc aha”.} Hai hiện tượng xuất hiện mà không hề được lập trình sẵn:
\begin{itemize}
	\item \textbf{Độ dài suy nghĩ tăng dần:} Độ dài trung bình của câu trả lời tăng đều đặn theo quá trình huấn luyện, từ vài trăm lên hàng nghìn token. Mô hình tự học được rằng \emph{suy nghĩ lâu hơn} giúp giải đúng nhiều hơn -- tức là tự học cách tận dụng tính toán lúc suy luận.
	\item \textbf{Hành vi phản tư:} Mô hình bắt đầu tự kiểm tra lại, quay lại xem xét các bước trước và thử cách tiếp cận khác. Nhóm tác giả ghi lại một “khoảnh khắc aha” khi một phiên bản trung gian của mô hình tự viết giữa chừng lời giải: \emph{“Wait, wait. Wait. That's an aha moment I can flag here.”} rồi đánh giá lại toàn bộ cách giải.
\end{itemize}

\begin{center}
\bookimage[0.98\textwidth]{r1_zero_response_length}{Hai biểu đồ: (a) độ chính xác trên AIME của DeepSeek-R1-Zero trong quá trình huấn luyện, gồm đường pass@1 và đường cons@16 (bỏ phiếu đa số trên 16 mẫu), cùng một đường ngang thể hiện mức của người tham gia thi; (b) độ dài trung bình mỗi câu trả lời tăng dần từ vài trăm lên hơn 10 nghìn token theo số bước huấn luyện.}
\captionof{figure}{DeepSeek-R1-Zero trong quá trình học tăng cường: độ chính xác AIME tăng và vượt mức người tham gia thi (trái), đồng thời độ dài câu trả lời tăng dần -- mô hình tự học cách “suy nghĩ lâu hơn” (phải). Nguồn: DeepSeek-AI~\cite{deepseekai2025r1}, phiên bản Nature, giấy phép CC BY 4.0.}
\label{fig:r1_zero_length}
\end{center}

\paragraph{Hạn chế.} Chuỗi suy nghĩ của R1-Zero \textbf{khó đọc} và thường \textbf{trộn lẫn nhiều ngôn ngữ} (ví dụ tiếng Anh và tiếng Trung trong cùng một đoạn).

\subsection{DeepSeek-R1: quy trình bốn giai đoạn}
\label{ssec:r1_pipeline}
Để có một mô hình vừa suy luận mạnh vừa thân thiện với người dùng, DeepSeek-R1 dùng quy trình đa giai đoạn:
\begin{enumerate}
	\item \textbf{Khởi động lạnh (cold start):} Tinh chỉnh DeepSeek-V3-Base trên vài nghìn ví dụ chuỗi suy luận dài, được định dạng dễ đọc (kèm phần tóm tắt cuối). Bước này tránh giai đoạn RL ban đầu thiếu ổn định và đưa “tiên nghiệm” của con người về cách trình bày suy luận vào mô hình.
	\item \textbf{RL hướng suy luận:} GRPO như R1-Zero trên toán, lập trình, khoa học, logic; thêm \textbf{phần thưởng nhất quán ngôn ngữ} (tỷ lệ từ thuộc ngôn ngữ mục tiêu trong chuỗi suy nghĩ) để giảm trộn lẫn ngôn ngữ -- chấp nhận giảm nhẹ hiệu năng để dễ đọc hơn.
	\item \textbf{Lấy mẫu loại bỏ + SFT:} Từ checkpoint RL, sinh nhiều lời giải và chỉ giữ lời giải đúng, dễ đọc: khoảng \textbf{600 nghìn mẫu suy luận}; cộng thêm khoảng \textbf{200 nghìn mẫu không phải suy luận} (viết, hỏi đáp, dịch\ldots) từ pipeline của DeepSeek-V3. Tinh chỉnh lại DeepSeek-V3-Base trên khoảng 800 nghìn mẫu này.
	\item \textbf{RL cho mọi tình huống:} Vòng RL thứ hai kết hợp phần thưởng dựa trên luật (cho suy luận) và reward model (cho tính hữu ích và vô hại của câu trả lời tổng quát). Tính hữu ích chỉ đánh giá phần tóm tắt cuối; tính vô hại đánh giá toàn bộ, kể cả chuỗi suy nghĩ.
\end{enumerate}
Kết quả: DeepSeek-R1 đạt \textbf{79.8\%} pass@1 trên AIME 2024 và \textbf{97.3\%} trên MATH-500, ngang OpenAI o1 (bản 12/2024).

\begin{center}
\bookimage[0.98\textwidth]{deepseek_r1_pipeline}{Sơ đồ khối quy trình huấn luyện DeepSeek-R1 theo bốn cột: (1) từ DeepSeek-V3-Base, học tăng cường với prompt suy luận và phần thưởng chính xác/định dạng tạo ra DeepSeek-R1-Zero, rồi lấy mẫu và lọc tạo dữ liệu khởi động lạnh; (2) SFT trên chuỗi suy luận dài tạo Dev1, sau đó RL với phần thưởng dựa trên luật và nhất quán ngôn ngữ tạo Dev2; (3) lấy mẫu dữ liệu suy luận và không suy luận từ Dev2 và DeepSeek-V3; (4) SFT rồi RL với phần thưởng dựa trên luật và phần thưởng theo sở thích để ra DeepSeek-R1.}
\captionof{figure}{Quy trình huấn luyện nhiều giai đoạn của DeepSeek-R1-Zero và DeepSeek-R1. Nguồn: DeepSeek-AI~\cite{deepseekai2025r1}, phiên bản Nature, giấy phép CC BY 4.0.}
\label{fig:r1_pipeline}
\end{center}

\subsection{Chưng cất so với RL chạy thẳng trên mô hình nhỏ}
\label{ssec:r1_distill}
Dùng chính 800 nghìn mẫu ở giai đoạn 3, nhóm tác giả tinh chỉnh (chỉ SFT, không RL) các mô hình Qwen2.5 và Llama 3 cỡ 1.5B, 7B, 8B, 14B, 32B và 70B. DeepSeek-R1-Distill-Qwen-32B đạt 72.6\% trên AIME 2024.

Thí nghiệm đối chứng quan trọng nhất của phần này: áp dụng RL quy mô lớn (hơn 10 nghìn bước) trực tiếp lên \emph{cùng} mô hình nền Qwen2.5-32B-Base chỉ đạt 47.0\% trên AIME 2024 -- kém bản chưng cất hơn 25 điểm, với chi phí tính toán lớn hơn nhiều lần. Đây không phải một quan sát vụn vặt; nó nói lên điều gì đó về cơ chế của hai phương pháp.

\paragraph{Vì sao khoảng cách lớn đến thế.}
Gradient của hai phương pháp bị chặn bởi hai đại lượng khác nhau. Với RLVR và phần thưởng nhị phân, đóng góp của một câu hỏi $q$ vào gradient tỉ lệ với mức bất đồng trong nhóm $G$ lời giải được lấy mẫu: nếu $k$ trong $G$ lời giải đúng thì độ lệch chuẩn phần thưởng của nhóm là
\[
\sigma_q \;=\; \frac{\sqrt{k\,(G-k)}}{G},
\]
bằng 0 đúng khi $k = 0$ hoặc $k = G$ (mục~\ref{ssec:grpo_variants} khai triển đẳng thức này). Một mô hình 32B chưa qua hậu huấn luyện suy luận có xác suất sinh ra lời giải đúng cho một bài AIME gần bằng 0, nên với hầu hết câu hỏi khó trong chương trình học, $k = 0$ và câu hỏi đó \emph{không đóng góp gì}. RL chỉ bơm thêm khối lượng xác suất cho những lời giải mà mô hình vốn đã thỉnh thoảng lấy mẫu trúng; nó không có kênh nào để tạo ra một lối giải mà mô hình chưa bao giờ sinh ra. Chi phí sinh $G$ chuỗi dài hàng nghìn token cho những câu hỏi ấy bị lãng phí hoàn toàn.

Cross-entropy trên trace của thầy không có cơ chế chặn này. Gradient $-\nabla_\theta \log \pi_\theta(y^{\text{thầy}}\mid q)$ khác 0 kể cả khi $\pi_\theta(y^{\text{thầy}}\mid q)$ nhỏ tùy ý -- thực tế càng nhỏ thì gradient càng lớn. Chưng cất vì vậy \emph{đưa vào} những khuôn mẫu lời giải nằm ngoài vùng mà trò tự lấy mẫu được, còn RL chỉ sắp xếp lại khối lượng xác suất bên trong vùng đó. Chi phí cũng chênh lệch theo cùng chiều: chưng cất trả tiền sinh dữ liệu từ thầy đúng một lần rồi huấn luyện bằng SFT thông thường, trong khi RL phải sinh $G$ chuỗi cho mỗi câu hỏi ở \emph{mỗi} bước cập nhật.

\paragraph{Hệ quả thực hành và giới hạn.}
Khi đã có sẵn một mô hình thầy mạnh và ngân sách cố định, chưng cất gần như luôn thắng RL chạy thẳng trên mô hình nhỏ. Nhưng trò bị chặn trên bởi thầy: chưng cất không tạo ra năng lực mà thầy không có, nên để đẩy biên hiện tại vẫn cần mô hình nền mạnh cộng RL quy mô lớn -- chính là quy trình tạo ra người thầy. Hai hướng cũng không loại trừ nhau. Tinh chỉnh trước để nâng $\pi_\theta$ lên đủ cao cho $k > 0$ trên nhiều câu hỏi hơn, rồi mới chạy RL, chính là lý do R1 có giai đoạn khởi động lạnh. Đọc theo cùng khung này, dynamic sampling của DAPO không chữa được vấn đề $k = 0$; nó chỉ ngừng trả tiền cho nó.

Mối liên hệ giữa chưng cất và học tăng cường được phân tích tiếp ở mục~\ref{sec:llm_distillation} và mục~\ref{sec:on_policy_distillation}; câu hỏi liệu học tăng cường có mở rộng được năng lực hay không được bàn ở mục~\ref{ssec:rl_capability_debate}.

\paragraph{Những thử nghiệm không thành công.} Nhóm DeepSeek cũng chia sẻ thất bại: (1) \textbf{PRM} gặp các vấn đề đã nêu ở mục~\ref{sec:verifiers}; (2) \textbf{Monte Carlo Tree Search} kiểu AlphaGo khó mở rộng vì không gian sinh token lớn hơn nhiều so với cờ vây và mô hình giá trị chi tiết rất khó huấn luyện.

Các cải tiến thuật toán xuất hiện sau R1 -- và câu hỏi liệu chúng có đang mở rộng năng lực của mô hình hay không -- đã được trình bày ở mục~\ref{sec:rl_algorithms} và mục~\ref{sec:rl_limits}.

\bigskip
\hrule
\bigskip

\begin{center}
	\textbf{\Large KẾT THÚC CHƯƠNG \thechapter}
\end{center}

\textit{Chương này đi từ các verifier kiểm chứng lời giải (ORM, PRM), qua các quy luật mở rộng tính toán lúc suy luận, đến các phương pháp tự học suy luận (STaR, ReST$^{EM}$, Quiet-STaR), rồi khép lại bằng DeepSeek-R1 -- trường hợp đầu tiên mà toàn bộ công thức huấn luyện một mô hình suy luận được mô tả công khai. Điểm chung của cả chương là một trục mở rộng thứ hai bên cạnh số tham số: lượng tính toán mà mô hình được phép tiêu lúc trả lời. Nhưng dù suy luận giỏi đến đâu, mô hình vẫn chỉ biết những gì đã nằm trong trọng số của nó; đưa tri thức cập nhật từ bên ngoài vào là chủ đề của chương tiếp theo về tìm kiếm và sinh tăng cường.}
